\section{Background}
\label{sec:background}

\subsection{Agent workflow graphs}
Many agent frameworks represent behavior as a directed graph, sometimes called
a \emph{state graph} or \emph{workflow graph}. Nodes correspond to steps such as:
an LLM call, a tool invocation, a router/selector that branches based on state,
or a human approval node. Edges correspond to transitions between these steps,
including conditional branches, parallel fan-outs, or loop back-edges.

This explicit control-flow model is a key difference between graph-orchestrated
agents and monolithic ``single prompt'' LLM applications: verification can be
performed over the topology without running the agent.

\subsection{Structural verification}
Structural properties are those that depend only on graph topology and type
annotations. Typical examples include:
(i) reachability (can a node be reached from the entry point?),
(ii) dead ends (nodes with no outgoing transitions),
(iii) isolation (must a sensitive tool be preceded by a gate node?), and
(iv) sanity constraints on router nodes.

These checks are inexpensive: they can be implemented with standard graph
algorithms (BFS/DFS) and run in time linear in the number of nodes and edges.

\subsection{Temporal monitors}
Temporal logics such as linear temporal logic (LTL) provide a formal language to
specify constraints over execution traces \citep{pnueli1977temporal}. A common
verification route is to compile temporal constraints into automata and then
reason over traces or over the product of a system model and the automaton.

Agentproof implements a lightweight monitoring path: a small temporal policy DSL
is compiled into a deterministic finite automaton (DFA) and evaluated over an
event stream representing tool calls and other agent actions. This design is
inspired by automata-theoretic monitoring approaches (e.g., \citep{vardi1994automata})
but intentionally focuses on a small, practical fragment.

\subsection{Related work (brief)}
Runtime guardrails, policy engines, and tool-call interceptors enforce safety
online, often with logging and audit support. Agentproof is complementary: it
targets \emph{static} properties of the workflow design that can be checked
before deployment, reducing the chance that unsafe topologies ship in the first
place.

